\section{Introduction}
\label{sec:intro}

NIST has standardized ML-KEM~\cite{fips203} and ML-DSA~\cite{fips204} for
post-quantum cryptography. They are entering widely used libraries and must run in
\emph{constant time}: if secrets influence branches, addresses, or instruction
latency, timing and microarchitectural observations can recover them~\cite{kocher}.
Rejection sampling, modular reduction, and polynomial arithmetic make such
mistakes easy in PQC.

\paragraph{The hook: KyberSlash.}
KyberSlash~\cite{kyberslash} is a family of timing leaks from a secret-dependent
integer division in several Kyber and ML-KEM implementations~\cite{kyber}. Control
flow and access pattern are identical for every input; the leak lives in the
\emph{latency} of one division instruction, which on many platforms depends on its
operands. A classic Valgrind/Memcheck taint screen~\cite{valgrind,ctgrind} flags
secret data reaching a branch or address, but KyberSlash slips past it: the code is
structurally clean while leaking through data-dependent timing.

\paragraph{The problem this creates for tooling.}
KyberSlash exposes three tooling problems. First, single checks have blind spots:
structural taint misses variable-latency instructions, while statistical timing
does not localize the cause. Second, verdicts are build-sensitive: optimization can
remove a debug-only branch or inline a leak into the wrong frame. Third, stacking
alarms tempts over-claiming: many PQC ``failures'' are variable-time by design,
such as ML-DSA rejection sampling~\cite{fips204,dilithium}.

\paragraph{Our approach.}
We present CT-KAT, a fail-closed constant-time screening framework for PQC. Under
one YAML configuration it binds three complementary checkers---a Valgrind/Memcheck
structural taint check, a warn-only assembly scan for variable-latency instruction
candidates, and a dudect-style statistical timing test~\cite{dudect}---and indexes
every structural verdict by a compiler $\times$ cflags build matrix. A failure is
never auto-labeled a key-recovery leak: leak sites are checked against a cited
accepted-variable-time registry, and anything unrecognized remains
\emph{needs-analysis} until a human rules (Figure~\ref{fig:pipeline}).

The phenomena are known: KyberSlash established variable-time division as a leak
class~\cite{kyberslash}, build-sensitivity is folklore, and dudect~\cite{dudect},
ctverif~\cite{ctverif}, and Binsec/Rel~\cite{binsecrel} predate this work. Our
contribution is their \emph{integration}, the \emph{triage discipline} that prevents
over-claiming, and its self-validation on PQClean~\cite{pqclean} targets---most
sharply, the default-deny taxonomy's refusal to over-claim a leak on ML-DSA-65
(Section~\ref{sec:results-coverage}).

\paragraph{Contributions.}
\begin{enumerate}
  \item \textbf{An integrated, build-matrix-indexed framework} binding a structural
    taint check, a variable-latency assembly scan, and a statistical timing test
    across a compiler $\times$ cflags matrix under one YAML config.
  \item \textbf{A default-deny triage taxonomy with a cited registry:} unrecognized
    leak sites stay at needs-analysis, making over-claiming structurally impossible.
  \item \textbf{A validated PQC corpus with the ML-DSA self-validation:} Docker runs
    over three families and five verdict classes (Table~\ref{tab:corpus}), a
    coverage analysis (Table~\ref{tab:coverage}), and a dudect appendix
    (Table~\ref{tab:dudect}).
\end{enumerate}
